\paragraph{Theorem D.2.30 (Proved: every finite additive lattice and every free twist).}
The fixed-$\gamma$ restriction in D.2.29 can also be removed.
All $102$ marked finite additive positions on the $172$ natural
candidate rows have the following integral description.
Let $d$ be the finite monodromy order, $q$ its index from
(D.HY2), and put
\[
 \gamma^*(\operatorname{Fix}A)=g\mathbb Z,\qquad
 g\in\{1,q\},\qquad q_v=q/g.
                                                               \tag{D.AF1}
\]
Here $g$ denotes the positive generator of that image.
There are $68$ positions with $g=1$ and $34$ with $g=q>1$.
Choose integral bases of the fixed and coinvariant lattices
so that
\[
\begin{gathered}
 \operatorname{Fix}A=\mathbb Z\delta\oplus\mathbb Z\ell,
 \qquad \gamma^*(\ell)=g,\\
 \Lambda/(A-1)\Lambda=\mathbb Z\beta\oplus\mathbb Z\eta,
 \qquad [\delta]=q_v\beta,\quad [\ell]=g\eta,\\
 N_m\beta=(m/q_v)\delta,
 \qquad N_m\eta=(m/g)\ell.
\end{gathered}                                                \tag{D.AF2}
\]
The lifts $\beta,\eta$ are part of the retained integral
marking. In particular $\gamma^*(\eta)=1$.

For $m=dk>0$, every power vector is $v=r\delta+n\ell$.
Put
\[
 t=\gcd(m,q_vr,gn).
\]
The exact free-action criterion is
\[
 \gcd(m,n,r)=1,\qquad d\mid m/t.                            \tag{D.AF3}
\]
The finite quotient and its normal line have precisely the
groups and the order-$t$ Chern class in (D.UF2).
Their full normal Thom kernel, on the fibre generator $s$
and the sweeps $D,L'$ of $\delta,\ell$, has the three columns
\[
 \begin{pmatrix}bgn/q_v\\m/q_v\\0\end{pmatrix},\qquad
 \begin{pmatrix}-br\\0\\m/g\end{pmatrix},\qquad
 \begin{pmatrix}0\\r\\n\end{pmatrix}.                     \tag{D.AF4}
\]
All entries are integers. The actual Seifert numerator is
$gn$, rather than $n$.

Consequently all $198$ minimal original/log additive-filling
patterns on all $148$ nonstar rows of the natural candidate
catalogue have a complete integral homology algorithm, for
every logarithmic multiplicity and every free affine twist.
Use (D.AF4) at each finite site, the full regular-log columns
of D.2.25 at arbitrarily many regular sites, and the single
relation (D.UF4). Take the torsion $T$ of the resulting Smith
cokernel, and take $r,A_1$ from the Seifert matrix with the
actual numerators $g_i n_i$. The answer is exactly (D.UF5),
including its zero-rational-Euler case. The complete
fundamental group remains D.2.23. The data of all $198$
patterns and all $102$ integral finite lattices are saved in
\path{records/manifolds_all_finite_logs.json}.

For one finite log and no regular logs there is also a closed
Smith formula, retaining an arbitrary integer clutch $e$.
Set
\[
 n'=n-(m/g)e,
 \qquad h=\gcd(b,m/g,n').                                  \tag{D.AF5}
\]
If $n'\ne0$, then
\[
\begin{gathered}
 \pi_1(X)=\mathbb Z/(g|n'|),\qquad
 H_1(X)=H_4(X)=\mathbb Z/(g|n'|),\\
 H_2(X)=H_3(X)=
 \mathbb Z/h\oplus\mathbb Z/\bigl(bg|n'|t/(qh)\bigr),
 \qquad H_5(X)=0.
\end{gathered}                                               \tag{D.AF6}
\]
When $n'=0$, instead
\[
\begin{gathered}
 \pi_1(X)=\mathbb Z,\qquad H_1(X)=H_5(X)=\mathbb Z,\\
 H_2(X)=H_3(X)=\mathbb Z^2\oplus
                         \mathbb Z/\gcd(b,m/g),\qquad
 H_4(X)=\mathbb Z^2.
\end{gathered}                                               \tag{D.AF7}
\]
In both displays $H_0=H_6=\mathbb Z$, and a cyclic summand of
order one is omitted. Every case with $g>1$ is necessarily
nonsimply connected, so this extension does not alter D.2.28.

\begin{proof}
The $g=1$ basis and its full marked comparison were proved in
D.2.26 and D.2.29. For the remaining positions the exact
matrices give
\[
 g=q,\qquad
 \Lambda/(A-1)\Lambda=\mathbb Z[\delta]\oplus\mathbb Z[\gamma],
 \qquad
 \ell=(q/d)N_d\gamma\in\Lambda.                            \tag{D.AF8}
\]
Thus take $\beta=\delta$ and $\eta=\gamma$.
The vectors $\delta,\ell$ are an integral basis of the fixed
lattice, and their coinvariant images are $\delta,q\gamma$.
These assertions are checked by unimodular Smith operations
for all $34$ actual local matrices, not only by rational
conjugacy. They prove (D.AF2). In all $102$ positions the
nonzero Smith factors of both $A-1$ and
$\bigwedge^2A-1$ are units. Hence the boundary groups used
in D.2.29 are again free of the stated ranks.

The invariant exterior form satisfies
\[
 -\omega(\beta,\ell)=bg/q_v,\qquad
 \omega(\eta,\delta)=b,\qquad
 \omega(\eta,\ell)=0.
\]
For $g=q$ these follow directly from (D.AF8), invariance of
$\omega$, and $\omega(\delta,\ell)=-bq$.
For $g=1$ they are the earlier identities. The required
divisibility $q_v\mid b$ is part of the same integral
catalogue, so the first coefficient of (D.AF4) is integral.

The fixed quotient torus has lattice $P\Lambda$, and the
coordinates of $v$ in it are $(q_vr,gn)$ by (D.AF2).
Repeating the actual-power fixed-point argument of D.2.29
therefore gives (D.AF3). The integer $t$ divides $q$ under
this freeness condition. The quotient surface has
\[
 H_1(F)=\langle\beta,\eta,c\mid
                   mc-q_vr\beta-gn\eta=0\rangle
       =\mathbb Z^2\oplus\mathbb Z/t.
\]
Its normal root character has value $1/t$ on the torsion
generator, exactly as before. Thus the normal line has Chern
class of order $t$, and duality gives all the cap groups.

The three full Thom generators in (D.UF8), with $\eta$ in
place of the former $\ell$ coinvariant generator, now have
normal projections $(m/q_v)\delta,(m/g)\ell,v$.
Evaluation of their actual fibre parallelograms by $\omega$
gives (D.AF4). Their primitive relation is
\[
 (m/t)K_c-(q_vr/t)K_\beta-(gn/t)K_\eta=0,
\]
and their normal-projection index is still $(m/q)t$.
The boundary-to-cap map in degree two remains onto, including
the torsion summand, because the normal Chern class is
torsion. Thus the full Mayer--Vietoris and duality argument of
D.2.29 applies with the actual numerator $gn$ at every site.
This proves the all-order matrix algorithm.

To check its scope, each finite additive site can now be
log-transformed, including those with $g>1$; for example a
power along the primitive fixed quotient direction gives a
free action. An original cap is an additional choice exactly
when its line multidegree is zero and its component action
has one orbit. Those are precisely the original caps whose
full integral comparison is already proved in D.2.26.
Enumerating these choices gives $198$ patterns on all $148$
nonstar initial rows. It introduces no further multiplicity
restriction. Zero-free multiplicative sites and the unique
rank-two zero retain the marked minimal models already used
in that catalogue.

For the one-log formula eliminate $D=-be\,s$ in (D.AF4),
and subtract $e$ times the second column from the third.
On $(s,L')$ the matrix becomes
\[
 \begin{pmatrix}
  bgn'/q_v&-br&0\\
  0&m/g&n'
 \end{pmatrix}.                                            \tag{D.AF9}
\]
Changing $n$ by $(m/g)e$ is exactly a peripheral lift change
by the coinvariant $\eta$. Thus (D.AF3) is preserved and
$t=\gcd(m,q_vr,gn')$. If $n'\ne0$, the gcd of the entries is
$h=\gcd(b,m/g,n')$: any prime dividing both $m/g$ and $n'$
cannot divide $r$ by freeness. The gcd of the two-by-two
minors is
\[
 bg|n'|t/q.
\]
These are the two determinantal divisors and prove the
middle groups in (D.AF6). The Seifert presentation gives the
cyclic group of order $g|n'|$. If $n'=0$, the sole remaining
relation is $-brs+(m/g)L'=0$, whose torsion divisor is
$\gcd(b,m/g)$ since $r$ is prime to $m$ in this free case.
Its free rank one acquires the one free degree-one bridge
class, giving the rank two in (D.AF7). Duality gives the
remaining groups in both cases.

Finally, if $g=q>1$, then $q$ divides both $m$ and its actual
numerator $gn$. Such a multiple fibre cannot occur in a
simply connected Seifert output by D.2.23. This verifies
that the new sector consists entirely of nonsimply-connected
members, without assuming that the group is abelian when
there are several logarithmic sites.
\end{proof}

The exact program \path{code/manifolds_all_finite_logs.py}
checks all $34$ additional lattice bases, compares the
freeness criterion with $38675$ actual-matrix power tests
for $1\le k\le6$ and all residues, and compares $175$ full
unreduced normal/graph Smith matrices to (D.AF4), including
$35$ torsion-cap instances. It also checks $725$ one-log
Smith forms with five different integer clutch values against
(D.AF6)--(D.AF7). The finite checks certify the explicit
constant matrices and examples; the proof covers all orders
and all free integer twists. Initial weights outside the
finite sphere-candidate catalogue and nonminimal original
caps remain separate parts of the broader D.2 request.
